% chapter.tex (Chapter 05)
\chapter{প্রোগ্রামিং ভাষা}

\begin{learningbox}
এই অধ্যায় শেষে শিক্ষার্থী পারবে:
\begin{itemize}
\item প্রোগ্রাম, প্রোগ্রামিং ভাষা, ভাষার প্রজন্ম ও translator program ব্যাখ্যা করতে।
\item algorithm, flowchart, program development step ও programming model বুঝতে।
\item C program-এর structure, data type, constant, variable, operator ও expression ব্যবহার করতে।
\item input/output, condition, loop, array ও function দিয়ে basic C program লিখতে।
\item board-style code output, error correction, algorithm ও flowchart প্রশ্ন সমাধান করতে।
\end{itemize}
\end{learningbox}

\begin{center}
\fbox{\parbox{0.92\textwidth}{
\centering
\textbf{অধ্যায় মানচিত্র}\\[4pt]
Program concept $\rightarrow$ Language generation $\rightarrow$ Translator $\rightarrow$ Algorithm/Flowchart $\rightarrow$ Programming model $\rightarrow$ C basics $\rightarrow$ Control/Array/Function $\rightarrow$ Practical programs
}}
\end{center}

\begin{keypointbox}{High-yield অংশ}
\begin{itemize}
\item Machine, assembly, high-level, 4GL এবং compiler/interpreter/assembler তুলনা।
\item Algorithm ও flowchart symbol।
\item C program structure, data type, variable, operator, keyword।
\item \texttt{printf}, \texttt{scanf}, \texttt{if}, \texttt{switch}, \texttt{for}, \texttt{while}, array, function।
\item CQ-তে algorithm + flowchart + C code + output একসাথে চাইতে পারে।
\end{itemize}
\end{keypointbox}

% Chapter 05 overview
\section{প্রোগ্রামিং ভাষা — পরিচিতি}

\begin{learningbox}
এই overview শেষে শিক্ষার্থী Chapter 5-এর scope বুঝবে: program, programming language, language generation, translator, algorithm, flowchart এবং C programming।
\end{learningbox}

\begin{center}
\fbox{\parbox{0.9\textwidth}{
\centering
\textbf{Chapter flow placeholder}\\[4pt]
Problem $\rightarrow$ Algorithm $\rightarrow$ Flowchart $\rightarrow$ C Code $\rightarrow$ Compile/Run $\rightarrow$ Output
}}
\end{center}

\begin{keypointbox}{এই অধ্যায়ে বেশি গুরুত্ব দাও}
\begin{itemize}
\item ভাষার প্রজন্ম ও translator program।
\item algorithm ও flowchart symbol।
\item C program-এর গঠন ও basic syntax।
\item condition, loop, array, function।
\item practical program: যোগফল, বড় সংখ্যা, জোড়-বিজোড়, factorial, series, array sum।
\end{itemize}
\end{keypointbox}

% ১. প্রোগ্রামের ধারণা
\section{প্রোগ্রামের ধারণা}

\begin{learningbox}
এই টপিক শেষে শিক্ষার্থী program, instruction, source/object/executable program এবং software-এর সম্পর্ক ব্যাখ্যা করতে পারবে।
\end{learningbox}

\begin{definitionbox}{প্রোগ্রাম}
কোনো নির্দিষ্ট কাজ করানোর জন্য কম্পিউটারকে দেওয়া ধারাবাহিক নির্দেশসমষ্টিকে প্রোগ্রাম বলা হয়।
\end{definitionbox}

\begin{center}
\begin{tabular}{|p{0.24\textwidth}|p{0.34\textwidth}|p{0.28\textwidth}|}
\hline
\textbf{ধরন} & \textbf{অর্থ} & \textbf{উদাহরণ} \\
\hline
Source program & programmer যে code লেখে & \texttt{sum.c} \\
\hline
Object program & translated machine-level file & object code \\
\hline
Executable program & সরাসরি run করা যায় & \texttt{sum.exe} \\
\hline
\end{tabular}
\end{center}

\begin{examplebox}{বাস্তব উদাহরণ}
Calculator app-এ user দুইটি সংখ্যা দিলে যোগফল দেখায়। input নেওয়া, যোগ করা এবং output দেখানোর নির্দেশগুলোই program।
\end{examplebox}

\begin{practicebox}
\begin{enumerate}
\item Program কী?
\item Source program ও object program-এর পার্থক্য লেখ।
\item Program ও software-এর সম্পর্ক ব্যাখ্যা কর।
\end{enumerate}
\end{practicebox}

% ২. প্রোগ্রামিং ভাষার ধারণা
\section{প্রোগ্রামিং ভাষার ধারণা}

\subsection{প্রোগ্রামিং ভাষা কী}
\begin{itemize}
\item মানব ভাষা বনাম কম্পিউটারের ভাষা
\item কেন কম্পিউটার সরাসরি বাংলা/ইংরেজি বুঝে না
\item Binary language ও instruction writing
\end{itemize}

\subsection{Program writing rules এবং প্রয়োজনীয়তা}
\begin{itemize}
\item লেখা ও সিনট্যাক্স নিয়ম
\item programming language-এর প্রয়োজনীয়তা
\end{itemize}

% ৩. বিভিন্ন প্রজন্মের প্রোগ্রামিং ভাষা
\section{বিভিন্ন প্রজন্মের প্রোগ্রামিং ভাষা}

\subsection{প্রথম প্রজন্ম — 1GL}
\begin{itemize}
\item Machine Language (0 ও 1)
\item Machine dependent, লেখা কঠিন
\end{itemize}

\subsection{দ্বিতীয় প্রজন্ম — 2GL}
\begin{itemize}
\item Assembly Language, mnemonic code
\item Assembler দিয়ে রূপান্তর
\end{itemize}

\subsection{তৃতীয় প্রজন্ম — 3GL}
\begin{itemize}
\item High Level Language: C, C++, Java, Fortran, BASIC ইত্যাদি
\item Compiler/Interpreter প্রয়োজন
\end{itemize}

\subsection{চতুর্থ প্রজন্ম — 4GL}
\begin{itemize}
\item Very high level: SQL, report/query tools
\end{itemize}

\subsection{পঞ্চম প্রজন্ম — 5GL}
\begin{itemize}
\item Natural language / AI ভিত্তিক ভাষা, logic programming
\end{itemize}

% ৪. নিম্নস্তরের ভাষা
\section{নিম্নস্তরের ভাষা}

\begin{learningbox}
এই টপিক শেষে শিক্ষার্থী low-level language-এর ধরন, সুবিধা ও সীমাবদ্ধতা লিখতে পারবে।
\end{learningbox}

\begin{definitionbox}{নিম্নস্তরের ভাষা}
যে programming language hardware বা machine-এর কাছাকাছি এবং machine dependent, তাকে নিম্নস্তরের ভাষা বলা হয়। Machine language ও assembly language এর অন্তর্ভুক্ত।
\end{definitionbox}

\begin{center}
\fbox{\parbox{0.88\textwidth}{
\centering
\textbf{Tree placeholder}\\[4pt]
Low Level Language $\rightarrow$ Machine Language, Assembly Language
}}
\end{center}

\begin{keypointbox}{মনে রাখো}
\begin{itemize}
\item execution দ্রুত হতে পারে।
\item hardware control বেশি।
\item লেখা, বুঝা ও debug করা কঠিন।
\item machine পরিবর্তন হলে program বদলাতে হতে পারে।
\end{itemize}
\end{keypointbox}

% ৪.১ Machine Language
\section{Machine Language}

\begin{learningbox}
এই টপিক শেষে শিক্ষার্থী machine language-এর সংজ্ঞা, বৈশিষ্ট্য, সুবিধা ও অসুবিধা ব্যাখ্যা করতে পারবে।
\end{learningbox}

\begin{definitionbox}{Machine language}
কম্পিউটারের processor সরাসরি যে binary instruction বা 0 ও 1-এর ভাষা বুঝতে পারে, তাকে machine language বলা হয়।
\end{definitionbox}

\begin{examplebox}{ধারণা}
\texttt{10110000 01100001} ধরনের binary pattern processor-এর কাছে instruction/data হিসেবে কাজ করতে পারে। মানুষ এটি সহজে বুঝতে পারে না, কিন্তু machine সরাসরি বুঝতে পারে।
\end{examplebox}

\begin{center}
\begin{tabular}{|p{0.30\textwidth}|p{0.50\textwidth}|}
\hline
\textbf{সুবিধা} & দ্রুত execute হয়; translator লাগে না। \\
\hline
\textbf{অসুবিধা} & লেখা, মনে রাখা, debug ও maintain করা কঠিন; machine dependent। \\
\hline
\end{tabular}
\end{center}

\begin{mistakebox}{সাধারণ ভুল}
Machine language দ্রুত হলেও programmer-friendly নয়। দ্রুত execution আর সহজ programming একই জিনিস নয়।
\end{mistakebox}

% ৪.২ Assembly Language
\section{Assembly Language}

\begin{learningbox}
এই টপিক শেষে শিক্ষার্থী assembly language, mnemonic code ও assembler-এর ভূমিকা ব্যাখ্যা করতে পারবে।
\end{learningbox}

\begin{definitionbox}{Assembly language}
যে low-level language-এ binary instruction-এর বদলে mnemonic code যেমন \texttt{ADD}, \texttt{MOV}, \texttt{SUB} ব্যবহার করা হয়, তাকে assembly language বলা হয়।
\end{definitionbox}

\begin{center}
\fbox{\parbox{0.88\textwidth}{
\centering
\textbf{Translation placeholder}\\[4pt]
Assembly source $\rightarrow$ Assembler $\rightarrow$ Machine code/Object code
}}
\end{center}

\begin{center}
\begin{tabular}{|p{0.26\textwidth}|p{0.30\textwidth}|p{0.30\textwidth}|}
\hline
\textbf{ভিত্তি} & \textbf{Machine} & \textbf{Assembly} \\
\hline
রূপ & 0 ও 1 & mnemonic code \\
\hline
Translator & লাগে না & assembler লাগে \\
\hline
বোঝা & খুব কঠিন & তুলনামূলক সহজ \\
\hline
\end{tabular}
\end{center}

% ৫. মধ্যস্তরের ভাষা
\section{মধ্যস্তরের ভাষা}

\begin{learningbox}
এই টপিক শেষে শিক্ষার্থী mid-level language কেন বলা হয় এবং C ভাষার ভূমিকা ব্যাখ্যা করতে পারবে।
\end{learningbox}

\begin{definitionbox}{মধ্যস্তরের ভাষা}
যে ভাষায় high-level language-এর মতো সহজ code লেখা যায়, আবার low-level operation বা hardware control-ও করা যায়, তাকে mid-level language বলা হয়।
\end{definitionbox}

\begin{examplebox}{C কেন mid-level}
C ভাষায় function, loop, data type দিয়ে structured program লেখা যায়; আবার pointer, bitwise operator ও memory-level কাজও করা যায়। তাই C-কে অনেক সময় mid-level language বলা হয়।
\end{examplebox}

\begin{boardbox}{বোর্ড ফোকাস}
C ভাষাকে mid-level বলার কারণ লিখতে high-level readability এবং low-level hardware control—দুই দিকই উল্লেখ করো।
\end{boardbox}

% ৬. উচ্চস্তরের ভাষা
\section{উচ্চস্তরের ভাষা}

\begin{learningbox}
এই টপিক শেষে শিক্ষার্থী high-level language-এর বৈশিষ্ট্য, সুবিধা ও translator প্রয়োজনীয়তা ব্যাখ্যা করতে পারবে।
\end{learningbox}

\begin{definitionbox}{উচ্চস্তরের ভাষা}
মানুষের ভাষা ও গণিতের notation-এর কাছাকাছি যে programming language, তাকে high-level language বলা হয়। এটি সাধারণত machine independent এবং compiler/interpreter দিয়ে machine code-এ রূপান্তরিত হয়।
\end{definitionbox}

\begin{keypointbox}{সুবিধা}
\begin{itemize}
\item লেখা ও বোঝা সহজ।
\item program development দ্রুত।
\item portability তুলনামূলক বেশি।
\item বড় application তৈরি করা সুবিধাজনক।
\end{itemize}
\end{keypointbox}

\begin{examplebox}{উদাহরণ}
C, C++, Java, Python, BASIC, Fortran, Visual Basic ইত্যাদি high-level language হিসেবে ব্যবহৃত হয়।
\end{examplebox}

% ৭. চতুর্থ প্রজন্মের ভাষা
\section{চতুর্থ প্রজন্মের ভাষা (4GL)}

\subsection{4GL কী}
\begin{itemize}
\item Very high level language, কম coding প্রয়োজন
\item Database, report, query ও business applications-এ ব্যবহার
\end{itemize}

\subsection{উদাহরণ ও ব্যবহার}
\begin{itemize}
\item SQL, Oracle environment, report generation, form design
\end{itemize}

% ৮. অনুবাদক প্রোগ্রাম
\section{অনুবাদক (Translator) প্রোগ্রাম}

\subsection{Translator Program কী}
\begin{itemize}
\item Source program-কে machine language বা object program-এ রূপান্তর
\item Compiler, Interpreter, Assembler সংক্ষেপে পরিচিতি
\end{itemize}

% ৮.১ Compiler
\section{Compiler}

\begin{learningbox}
এই টপিক শেষে শিক্ষার্থী compiler-এর কাজ, সুবিধা, অসুবিধা এবং interpreter-এর সাথে পার্থক্য লিখতে পারবে।
\end{learningbox}

\begin{definitionbox}{Compiler}
যে translator পুরো source program একসাথে machine code/object code-এ রূপান্তর করে এবং error list দেখায়, তাকে compiler বলা হয়।
\end{definitionbox}

\begin{center}
\fbox{\parbox{0.88\textwidth}{
\centering
\textbf{Compiler flow placeholder}\\[4pt]
Source code $\rightarrow$ Compiler $\rightarrow$ Object code $\rightarrow$ Linker $\rightarrow$ Executable
}}
\end{center}

\begin{keypointbox}{সুবিধা ও সীমাবদ্ধতা}
\begin{itemize}
\item একবার compile হলে executable দ্রুত run হয়।
\item object/executable file তৈরি হতে পারে।
\item সব error একসাথে দেখাতে পারে, তাই নতুনদের কাছে debugging কঠিন লাগতে পারে।
\end{itemize}
\end{keypointbox}

\begin{examplebox}{C language}
C program সাধারণত compiler দিয়ে machine code-এ রূপান্তরিত হয়, তারপর run করা হয়।
\end{examplebox}

% ৮.২ Interpreter
\section{Interpreter}

\begin{learningbox}
এই টপিক শেষে শিক্ষার্থী interpreter-এর কাজ এবং compiler থেকে পার্থক্য ব্যাখ্যা করতে পারবে।
\end{learningbox}

\begin{definitionbox}{Interpreter}
যে translator source program line by line translate ও execute করে এবং error পেলে সাধারণত সেই line-এ execution থামায়, তাকে interpreter বলা হয়।
\end{definitionbox}

\begin{center}
\begin{tabular}{|p{0.24\textwidth}|p{0.30\textwidth}|p{0.30\textwidth}|}
\hline
\textbf{ভিত্তি} & \textbf{Compiler} & \textbf{Interpreter} \\
\hline
Translation & পুরো program একসাথে & line by line \\
\hline
Execution speed & দ্রুত & তুলনামূলক ধীর \\
\hline
Error handling & error list একসাথে & line ধরে error দেখায় \\
\hline
Object code & তৈরি হতে পারে & সাধারণত তৈরি করে না \\
\hline
\end{tabular}
\end{center}

\begin{examplebox}{শেখার সুবিধা}
Interpreter line ধরে error দেখায় বলে beginners দ্রুত ভুল খুঁজতে পারে।
\end{examplebox}

% ৮.৩ Assembler
\section{Assembler}

\subsection{Assembler কী}
\begin{itemize}
\item Assembly language-কে machine language-এ রূপান্তর করে
\item mnemonic codes translate করে object program তৈরি হয়
\end{itemize}

% ৯. প্রোগ্রামের সংগঠন
\section{প্রোগ্রামের সংগঠন}

\begin{learningbox}
এই টপিক শেষে শিক্ষার্থী একটি C program-এর প্রধান অংশ এবং IPO model ব্যাখ্যা করতে পারবে।
\end{learningbox}

\begin{definitionbox}{Program organization}
একটি program-এর বিভিন্ন অংশ যেমন comment, header, declaration, main function, input, processing, output ও return statement সুশৃঙ্খলভাবে সাজানোর পদ্ধতিকে program organization বলা যায়।
\end{definitionbox}

\begin{center}
\fbox{\parbox{0.88\textwidth}{
\centering
\textbf{IPO model placeholder}\\[4pt]
Input $\rightarrow$ Processing $\rightarrow$ Output
}}
\end{center}

\begin{keypointbox}{C program-এর সাধারণ অংশ}
\begin{enumerate}
\item Documentation/comment
\item Preprocessor directive
\item Global declaration
\item \texttt{main()} function
\item Variable declaration
\item Input, processing, output
\item \texttt{return 0;}
\end{enumerate}
\end{keypointbox}

\begin{boardbox}{বোর্ড ফোকাস}
C program structure আঁকতে \texttt{\#include <stdio.h>} এবং \texttt{int main()} অংশ সাধারণত লিখতে হয়।
\end{boardbox}

% ১০. প্রোগ্রাম উন্নয়নের ধাপসমূহ
\section{প্রোগ্রাম উন্নয়নের ধাপসমূহ}

\begin{learningbox}
এই টপিক শেষে শিক্ষার্থী program development life cycle-এর ধাপগুলো ধারাবাহিকভাবে লিখতে পারবে।
\end{learningbox}

\begin{center}
\fbox{\parbox{0.92\textwidth}{
\centering
\textbf{Development flowchart placeholder}\\[4pt]
Problem definition $\rightarrow$ Analysis $\rightarrow$ Algorithm $\rightarrow$ Flowchart $\rightarrow$ Coding $\rightarrow$ Compile/Run $\rightarrow$ Testing/Debugging $\rightarrow$ Documentation/Maintenance
}}
\end{center}

\begin{keypointbox}{ধাপসমূহ}
\begin{itemize}
\item সমস্যা নির্ধারণ ও বিশ্লেষণ।
\item Input, process ও output নির্ধারণ।
\item Algorithm ও flowchart তৈরি।
\item Coding, compiling, testing ও debugging।
\item Documentation ও maintenance।
\end{itemize}
\end{keypointbox}

\begin{examplebox}{বাস্তব প্রয়োগ}
তিনটি সংখ্যার মধ্যে বড় সংখ্যা বের করার program লিখতে আগে input তিনটি সংখ্যা, process তুলনা, output বড় সংখ্যা—এভাবে IPO নির্ধারণ করলে code সহজ হয়।
\end{examplebox}

% ১১. অ্যালগরিদম
\section{অ্যালগরিদম}

\subsection{Algorithm-এর ধারণা}
\begin{itemize}
\item ধাপে ধাপে নির্দেশাবলী; start/stop, finite, logical steps
\end{itemize}

\subsection{Algorithm লেখার নিয়ম ও উদাহরণ}
\begin{itemize}
\item Step-1, Step-2 আকারে লেখা
\item উদাহরণ: দুই সংখ্যার যোগফল, বড় সংখ্যা নির্ণয়, factorial ইত্যাদি
\end{itemize}

% ১২. ফ্লোচার্ট
\section{ফ্লোচার্ট}

\subsection{Flowchart-এর ধারণা}
\begin{itemize}
\item Algorithm-এর চিত্রভিত্তিক প্রকাশ, symbol ব্যবহার
\end{itemize}

\subsection{Flowchart Symbol}
\begin{itemize}
\item Oval: Start/End, Parallelogram: Input/Output, Rectangle: Process
\item Diamond: Decision, Arrow: Flow, Circle: Connector
\end{itemize}

\subsection{Practice}
\begin{itemize}
\item যোগফল, বড় সংখ্যা, factorial, temperature conversion ইত্যাদি
\end{itemize}

% ১৩. প্রোগ্রাম ডিজাইন মডেল
\section{প্রোগ্রাম ডিজাইন মডেল}

\begin{learningbox}
এই টপিক শেষে শিক্ষার্থী structured, visual, object-oriented ও event-driven programming model-এর ধারণা তুলনা করতে পারবে।
\end{learningbox}

\begin{center}
\begin{tabular}{|p{0.24\textwidth}|p{0.42\textwidth}|p{0.22\textwidth}|}
\hline
\textbf{Model} & \textbf{ধারণা} & \textbf{উদাহরণ} \\
\hline
Structured & sequence, selection, loop; module-based & C \\
\hline
Visual & GUI control drag-and-drop করে design & Visual Basic \\
\hline
OOP & class ও object ভিত্তিক & C++, Java \\
\hline
Event-driven & user/system event অনুযায়ী কাজ & GUI app \\
\hline
\end{tabular}
\end{center}

\begin{center}
\fbox{\parbox{0.88\textwidth}{
\centering
\textbf{Model map placeholder}\\[4pt]
Problem type $\rightarrow$ Suitable programming model $\rightarrow$ Program design
}}
\end{center}

\begin{boardbox}{বোর্ড ফোকাস}
Programming model-এর প্রশ্নে শুধু নাম লিখবে না; কোন ধরনের problem/application-এ ব্যবহার হয় সেটিও লিখবে।
\end{boardbox}

% ১৩.১ Structured Programming
\section{Structured Programming}

\begin{itemize}
\item বড় problem-কে ছোট module-এ ভাগ করা
\item Control structures: sequence, selection, loop
\item Modular programming ও top-down approach
\end{itemize}

% ১৩.২ Visual Programming
\section{Visual Programming}

\begin{learningbox}
এই টপিক শেষে শিক্ষার্থী visual programming-এর GUI-based ধারণা ও সুবিধা লিখতে পারবে।
\end{learningbox}

\begin{definitionbox}{Visual programming}
যে programming environment-এ form, button, textbox ইত্যাদি graphical control ব্যবহার করে application design করা যায়, তাকে visual programming বলা হয়।
\end{definitionbox}

\begin{center}
\fbox{\parbox{0.88\textwidth}{
\centering
\textbf{GUI placeholder}\\[4pt]
Form window: Label, TextBox, Button, Output area
}}
\end{center}

\begin{keypointbox}{সুবিধা}
\begin{itemize}
\item interface দ্রুত design করা যায়।
\item drag-and-drop control ব্যবহার করা যায়।
\item event-based code লেখা সহজ হয়।
\end{itemize}
\end{keypointbox}

% ১৩.৩ Object Oriented Programming
\section{Object Oriented Programming}

\begin{itemize}
\item Object, Class, Method, Property
\item Encapsulation, Inheritance, Polymorphism
\item উদাহরণ: Java, C++
\end{itemize}

% ১৩.৪ Event Driven Programming
\section{Event Driven Programming}

\begin{itemize}
\item Event কী, event handler, GUI events (mouse, keyboard)
\item Event-driven application structure
\end{itemize}

% ১৪. C প্রোগ্রামিং ভাষা — প্রাথমিক ধারণা
\section{C প্রোগ্রামিং — প্রাথমিক ধারণা}

\begin{itemize}
\item C language কী, Dennis Ritchie
\item Mid level, structured, general-purpose language
\end{itemize}

% ১৪.১ C ভাষার বৈশিষ্ট্য
\section{C ভাষার বৈশিষ্ট্য}

\begin{itemize}
\item Mid level, structured, portable, fast, function-based
\item Low level operations possible, rich library, case sensitive
\end{itemize}

% ১৪.২ C Program Compiling
\section{C Program Compiling}

\begin{itemize}
\item Source code → Compiler → Object code → Linker → Executable
\item Error checking, run/execute, debugging ধারাবাহিকতা
\end{itemize}

% ১৪.৩ C Program-এর গঠন
\section{C Program-এর গঠন}

\begin{learningbox}
এই টপিক শেষে শিক্ষার্থী একটি basic C program-এর অংশ চিনতে ও লিখতে পারবে।
\end{learningbox}

\begin{verbatim}
#include <stdio.h>

int main()
{
    int a, b, sum;
    scanf("%d %d", &a, &b);
    sum = a + b;
    printf("Sum = %d", sum);
    return 0;
}
\end{verbatim}

\begin{center}
\begin{tabular}{|p{0.28\textwidth}|p{0.52\textwidth}|}
\hline
\textbf{অংশ} & \textbf{কাজ} \\
\hline
\texttt{\#include <stdio.h>} & standard input-output library যুক্ত করে \\
\hline
\texttt{int main()} & program execution শুরু হয় \\
\hline
Declaration & variable-এর type ও name ঘোষণা করে \\
\hline
\texttt{scanf}/\texttt{printf} & input নেওয়া ও output দেখানো \\
\hline
\texttt{return 0;} & program সফলভাবে শেষ বোঝায় \\
\hline
\end{tabular}
\end{center}

\begin{boardbox}{বোর্ড ফোকাস}
Code লেখার প্রশ্নে indentation, semicolon, brace এবং format specifier ঠিক রাখা জরুরি।
\end{boardbox}

% ১৪.৪ Data Type
\section{Data Type}

\begin{itemize}
\item Integer, Floating point, Character, Double, Void
\item Common types: int, float, double, char, void
\end{itemize}

% ১৪.৫ Constant
\section{Constant}

\begin{learningbox}
এই টপিক শেষে শিক্ষার্থী constant-এর ধরন ও ব্যবহার লিখতে পারবে।
\end{learningbox}

\begin{definitionbox}{Constant}
Program execution চলাকালে যার মান পরিবর্তন হয় না, তাকে constant বলা হয়।
\end{definitionbox}

\begin{center}
\begin{tabular}{|p{0.25\textwidth}|p{0.35\textwidth}|p{0.25\textwidth}|}
\hline
\textbf{ধরন} & \textbf{উদাহরণ} & \textbf{নোট} \\
\hline
Integer constant & \texttt{10}, \texttt{-5} & পূর্ণসংখ্যা \\
\hline
Floating constant & \texttt{3.14} & দশমিকসহ সংখ্যা \\
\hline
Character constant & \texttt{'A'} & single quote \\
\hline
String constant & \texttt{"ICT"} & double quote \\
\hline
Symbolic constant & \texttt{\#define PI 3.1416} & নাম দিয়ে মান \\
\hline
\end{tabular}
\end{center}

% ১৪.৬ Variable
\section{Variable}

\begin{learningbox}
এই টপিক শেষে শিক্ষার্থী variable declaration, initialization ও naming rule লিখতে পারবে।
\end{learningbox}

\begin{definitionbox}{Variable}
Program execution চলাকালে যার মান পরিবর্তন হতে পারে, এমন named memory location-কে variable বলা হয়।
\end{definitionbox}

\begin{examplebox}{Declaration ও initialization}
\texttt{int age;} declaration। \texttt{int age = 18;} declaration with initialization।
\end{examplebox}

\begin{keypointbox}{Naming rule}
\begin{itemize}
\item letter বা underscore দিয়ে শুরু হবে।
\item digit দিয়ে শুরু করা যাবে না।
\item keyword ব্যবহার করা যাবে না।
\item space থাকবে না।
\item C case sensitive।
\end{itemize}
\end{keypointbox}

\begin{mistakebox}{সাধারণ ভুল}
\texttt{2num}, \texttt{student mark}, \texttt{int} variable name হিসেবে valid নয়।
\end{mistakebox}

% ১৪.৭ Operator
\section{Operator}

\begin{learningbox}
এই টপিক শেষে শিক্ষার্থী C language-এর operator-এর ধরন ও ব্যবহার বুঝবে।
\end{learningbox}

\begin{definitionbox}{Operator}
Variable বা value-এর উপর নির্দিষ্ট operation সম্পাদনকারী symbol-কে operator বলা হয়।
\end{definitionbox}

\begin{center}
\begin{tabular}{|p{0.24\textwidth}|p{0.36\textwidth}|p{0.24\textwidth}|}
\hline
\textbf{ধরন} & \textbf{Operator} & \textbf{ব্যবহার} \\
\hline
Arithmetic & \texttt{+ - * / \%} & যোগ, বিয়োগ, ভাগ \\
\hline
Relational & \texttt{> < >= <= == !=} & তুলনা \\
\hline
Logical & \texttt{\&\& || !} & condition combine \\
\hline
Assignment & \texttt{=} & মান বসানো \\
\hline
Increment/Decrement & \texttt{++ --} & 1 বাড়ানো/কমানো \\
\hline
\end{tabular}
\end{center}

\begin{mistakebox}{সাধারণ ভুল}
\texttt{=} assignment operator, আর \texttt{==} equality test operator। condition-এ এদের গুলিয়ে ফেললে output ভুল হয়।
\end{mistakebox}

% ১৪.৮ Expression
\section{Expression}

\begin{itemize}
\item Expression কী, Arithmetic/Relational/Logical expressions
\item Operator precedence ও associativity
\end{itemize}

% ১৪.৯ Keyword
\section{Keyword}

\begin{itemize}
\item Keyword কী, reserved words, variable name হিসেবে ব্যবহার করা যাবে না
\item উদাহরণ: int, float, char, if, else, for, while, return, break
\end{itemize}

% ১৪.১০ Input Output Statement
\section{Input/Output Statement}

\begin{learningbox}
এই টপিক শেষে শিক্ষার্থী \texttt{printf} ও \texttt{scanf} ব্যবহার করে input-output করতে পারবে।
\end{learningbox}

\begin{definitionbox}{\texttt{printf}}
C language-এ output দেখানোর জন্য \texttt{printf()} function ব্যবহৃত হয়।
\end{definitionbox}

\begin{definitionbox}{\texttt{scanf}}
Keyboard থেকে input নেওয়ার জন্য \texttt{scanf()} function ব্যবহৃত হয়।
\end{definitionbox}

\begin{verbatim}
int n;
scanf("%d", &n);
printf("Number = %d", n);
\end{verbatim}

\begin{center}
\begin{tabular}{|p{0.18\textwidth}|p{0.30\textwidth}|p{0.30\textwidth}|}
\hline
\textbf{Specifier} & \textbf{Data type} & \textbf{Example} \\
\hline
\texttt{\%d} & int & \texttt{10} \\
\hline
\texttt{\%f} & float & \texttt{3.5} \\
\hline
\texttt{\%c} & char & \texttt{A} \\
\hline
\texttt{\%s} & string & \texttt{ICT} \\
\hline
\end{tabular}
\end{center}

\begin{mistakebox}{সাধারণ ভুল}
\texttt{scanf}-এ variable-এর আগে সাধারণত address operator \texttt{\&} দিতে হয়, যেমন \texttt{\&n}।
\end{mistakebox}

% ১৪.১১ Conditional Statement
\section{Conditional Statement}

\begin{itemize}
\item if, if...else, else if ladder, nested if
\item switch statement
\end{itemize}

\subsection{উদাহরণ}
\begin{verbatim}
if (mark >= 33) {
    printf("Pass");
} else {
    printf("Fail");
}
\end{verbatim}

% ১৪.১২ Loop Statement
\section{Loop Statement}

\begin{learningbox}
এই টপিক শেষে শিক্ষার্থী \texttt{for}, \texttt{while}, \texttt{do...while} loop-এর ব্যবহার লিখতে পারবে।
\end{learningbox}

\begin{definitionbox}{Loop}
একই statement বা statement block বারবার execute করার জন্য যে control structure ব্যবহৃত হয়, তাকে loop বলা হয়।
\end{definitionbox}

\begin{center}
\begin{tabular}{|p{0.20\textwidth}|p{0.44\textwidth}|p{0.24\textwidth}|}
\hline
\textbf{Loop} & \textbf{ধারণা} & \textbf{Condition check} \\
\hline
\texttt{for} & count জানা থাকলে সুবিধাজনক & আগে \\
\hline
\texttt{while} & condition true থাকা পর্যন্ত চলে & আগে \\
\hline
\texttt{do...while} & অন্তত একবার চলে & পরে \\
\hline
\end{tabular}
\end{center}

\begin{verbatim}
for (int i = 1; i <= 5; i++) {
    printf("%d ", i);
}
\end{verbatim}

\begin{mistakebox}{সাধারণ ভুল}
Loop condition update না করলে infinite loop হতে পারে। initialization, condition, increment/decrement খেয়াল রাখতে হবে।
\end{mistakebox}

% ১৪.১৩ Array
\section{Array}

\begin{itemize}
\item Array কী, index, one-dimensional, two-dimensional
\item declaration, initialization, input-output
\end{itemize}

\subsection{উদাহরণ}
\begin{verbatim}
int marks[5];
\end{verbatim}

% ১৪.১৪ Function
\section{Function}

\begin{itemize}
\item Function কী, library vs user-defined
\item declaration, definition, calling, parameter ও return value
\item সুবিধা: code reuse, modularity, debugging সহজ
\end{itemize}

% ১৫. Practical Program Practice
\section{Practical Program Practice}

\begin{itemize}
\item দুই সংখ্যার যোগফল, তিন সংখ্যার মধ্যে বড় সংখ্যা
\item জোড়-বিজোড় নির্ণয়, pass/fail নির্ণয়
\item for/while examples, factorial, multiplication table
\item array input-output, function-based program examples
\end{itemize}

% ১৬. অনুশীলনী
\section{অনুশীলনী}

\subsection{প্রশ্নের ধরন}
\begin{itemize}
\item MCQ, জ্ঞানমূলক, অনুধাবনমূলক, প্রয়োগমূলক, উচ্চতর দক্ষতামূলক
\item Algorithm লেখা, Flowchart আঁকা, C code লেখা ও ভুল code correction
\end{itemize}

\subsection{প্র্যাকটিস প্রোগ্রাম ও কার্যাবলী}
\begin{itemize}
\item Compiler/Interpreter পার্থক্য, ছোট C programs, output নির্ণয়
\end{itemize}


% Resource linked: resources/HSC ICT Chapter 05 Handnote.pdf
